\input{../preamble.tex}

\SetDate[10/04/2026]

\reversemarginpar
\setlength{\marginparsep}{.5cm}

\begin{document}
\pagestyle{fancy}
\fancyhead[L]{Tle STMG}
\fancyhead[C]{\textbf{Évaluation — Statistiques à deux variables}}
\fancyhead[R]{\today}

%\null\vspace{-30pt}
Consignes particulières : 
\begin{itemize}[label=$\bullet$]
	\item 
	La calculatrice est {autorisée}.
	\item
	L'évaluation fait \pageref{lastpage} page. La somme des points est \total{points}.
	\item
	Formules du coefficient directeur $a$ et de l'ordonnée à l'origine $b$ :
	\begin{align*}
		a = \dfrac{y_B-y_A}{x_B-x_A} && \et && b =y_A - ax_A.
	\end{align*}
\end{itemize}

\marginpar{[pts]}
\hrule



\exemulticols{7}{
	Deux ajustements affines sont proposés pour approximer le nuage de points ci-contre :
		\begin{align*}
			f(x) = x+1, && \et && g(x) = 0,5+1,2x.
		\end{align*}
	\begin{enumerate}
		\item
		Remplir les deux premières lignes du tableau, chaque colonne correspondant à un point du nuage de coordonnées $(x ; y)$.
		\item
		À quelle droite, $(d_1)$ ou $(d_2)$, correspond la fonction $f$ ? Justifier.
		\item
		Calculer l'erreur commise par chaque ajustement affine proposé à l'aide du tableau.
		\item
		Quel ajustement permet de mieux décrire le nuage de points ?
	\end{enumerate}
	\vfill
}{
	\null\vspace{-40pt}
	\begin{center}
	\begin{tikzpicture}
		\begin{axis}[xmin = 0, ymin=0, ymax=10, axis line style=->, extra y ticks={0}, extra x ticks={0}]
			\addplot[only marks, BLACK, mark size=2pt]  coordinates {
			(1,1)
			(1,2)
			(3,4)
			(4,4)
			(4,5)
			(6,8)
			(7,9)
			}; 
			
			\addplot[very thick, dashed, BLUE_E] expression[domain=0:7, samples=2]{1+x} node[right]{$(d_2)$};
			\addplot[very thick, dashed, RED_E] expression[domain=0:7, samples=2]{.5+1.2*x} node[right]{$(d_1)$};
		\end{axis}
	\end{tikzpicture}
	\end{center}
}{exe:l2-comparison}{
}
	\begin{center}
	
\def\arraystretch{1.3}
\setlength\tabcolsep{25pt}
	\begin{tabular}{|*{8}{c|}}\hline
		$x$ & & & & & & &
		\\\hline
		$y$ & & & & & & &
		\\\hline
		$\bigl[f(x)-y\bigr]^2$ & & & & & & &
		\\\hline
		$\bigl[g(x)-y\bigr]^2$ & & & & & & &
		\\\hline
	\end{tabular}
	\end{center}

\exe{4}{
	Pour chaque jeu de données ci-dessous, un ajustement affine est-il pertinent ? Justifier.
	
\newcommand{\scale}{.7}
	\begin{multicols}{2}
	\begin{center}
	\begin{tikzpicture}[scale=\scale]
		\begin{axis}[xmin = 0, xmax=13, ymin=700, ymax=1000, ticks = none, clip=false, axis line style=->,]
			\addplot[only marks, BLACK, mark size=2pt] table[x=anc, y=sal] {data/2-1.dat};
		\end{axis}
	\end{tikzpicture}
	
	(a)
	\end{center}
	
	\columnbreak
	
	\begin{center}
	\begin{tikzpicture}[scale=\scale]
		\begin{axis}[xmin = 0, xmax=10, ymin=970, ymax=1000, ticks = none, clip=false, axis line style=->,]
			\addplot[only marks, BLACK, mark size=2pt] table[x=anc, y=sal] {data/2-2.dat};
		\end{axis}
	\end{tikzpicture}
	
	(b)
	\end{center}
	\end{multicols}
}{exe:pertinentYN}{
}

\newpage

\exe{9}{
	Un médecin prescrit 500mg de paracétamol à sa patiente et surveille la quantité absorbée à l'aide de plusieurs prises de sang.
	À partir de l'heure 0, il reporte la quantité de médicament encore présente dans l'organisme après 1 ; 2 ; 5 ; 8 ; et 11 heures.

%	Leonardo Fibonacci élève des lapins de la race « Romarin » et étudie l'évolution de la population.
%	Partant de l'année 0, il recense chaque année le nombre de lapins et le note dans son carnet.
%	
%	Ne sachant relire son écriture, les données du tableau sont malheureusement incomplètes.
	 
	\begin{multicols}{2}
	
	\begin{center}
	\begin{tikzpicture}
	\begin{axis}[
 		axis line style=->,
		grid=both,
		minor y tick num=1,
		xtick distance=1,
		ytick distance=50,
		extra x ticks={0},
		extra y ticks={0},
		x=18pt,
		y=.4pt,
		xmin=0,
		xmax=11,
		ymin=0,
		ymax=500,
		xlabel = Heure,
		ylabel = Quantité de médicament (mg),
		ylabel near ticks,
		xlabel near ticks,
		clip=true,
		]
		\addplot[only marks, transparent, mark size=2pt] table[x=h, y=Mh] {data/2-402.dat};
	\end{axis}
	\end{tikzpicture}
	\end{center}
	
\def\arraystretch{1.25}
\setlength\tabcolsep{5pt}
	\hfill
	\begin{tabular}{|c|c|c|}
		\hline
		\textbf{Heure} & \textbf{\thead{Quantité encore \\ présente (mg)}} & \textbf{\thead{Logarithme de la \\ quantité encore présente}}
		\\\hline
		0 & 500 & \\ \hline
		1 & 300 & \\ \hline
		2 & 180 & \\\hline
		%3 & 108 & \\\hline
		%4 & 65 & \\\hline
		3 & & \\\hline
		4 & & \\\hline
		5 & 39 & \\\hline
		%6 & 23 & \\ \hline
		%7 & 14 & \\\hline
		6 & & \\ \hline
		7 & & \\\hline
		8 & 8 & \\\hline
		%9 & 5 & \\\hline
		%10 & 3 & \\\hline
		9 &  & \\\hline
		10 &  & \\\hline
		11 & 2 & \\\hline
	\end{tabular}
	\vfill\null
	\end{multicols}
	
	\begin{enumerate}
		\item
		Représenter les données du médecin dans le repère ci-dessus.
		\item
		Un ajustement affine des données vous semble-t-il pertinent ? Justifier.
		\item
		Suspectant une décroissance exponentielle, le médecin décide d'appliquer le logarithme à la quantité de médicament encore présente.
		Compléter le tableau en arrondissant les logarithmes à $10^{-2}$ près.
		\item
		Représenter le nouveau nuage de points dans le repère ci-dessous.
		\item 
		Déterminer l'équation d'une droite $(d)$ d'ajustement affine liant le logarithme de la quantité de médicament $\log(y)$ et l'heure $x$ sous la forme $\log(y)=ax+b$.
		Arrondir les paramètres $a$ et $b$ à $10^{-2}$.
		\item
		Remplir les données manquantes du tableau à l'aide de votre ajustement affine.
		Arrondir les logarithmes à $10^{-2}$ et les quantités de médicament à l'unité.
		\item
		Quelle quantité de médicament sera encore présente dans l'organisme après 24 heures ?
		Prédire la valeur à l'aide de votre ajustement affine.
	\end{enumerate}
	
	
	\begin{center}
	\begin{tikzpicture}
	\begin{axis}[
 		axis line style=->,
		grid=both,
		minor y tick num=1,
		xtick distance=2,
		ytick distance=1,
		%extra x ticks={0},
		extra y ticks={0},
		x=19pt,
		%y=5pt,
		xmin=0,
		xmax=24,
		ymin=-3,
		ymax=3,
		xlabel = Heure,
		ylabel = \thead{Logarithme de la \\ quantité de médicament},
		ylabel near ticks,
		%xlabel near ticks,
		clip=true,
		]
		\addplot[only marks, transparent, mark size=2pt] table[x=h, y=log] {data/2-402.dat};
		
		\addplot[thick, dashed, transparent] expression[domain=0:24, samples=2]{log10(500)+log10(.6)*x};
	\end{axis}
	\end{tikzpicture}
	\end{center}
	
	
}{exe:LuLu}{
}



%%%%%%%%%%%%

\label{lastpage}
\newpage
\fancyhead[C]{\textbf{Solutions}}
\shipoutAnswer

\end{document}
